\chapter{Discussion}\label{ch:discussion}

This chapter discusses the results in the same objective order used in Chapters~\ref{ch:methods} and~\ref{ch:results}. The discussion first considers the prototype as an integrated wearable platform, then the detection task, then the Shimmer/XIAO sensing evidence, and finally the embedded reminder feasibility.

\section{Objective 1: Prototype Feasibility and Engineering Trade-Offs}

The prototype addresses Objective~1 by integrating the main technical components required for a pelvis-focused reminder platform: local IMU sensing, BLE streaming, labelled Android collection, PC inspection, compact embedded classification, and local haptic actuation. This positions the work between sedentary-behaviour interventions \cite{stephenson2017technology_sedentary_meta} and wearable movement sensing \cite{papi2017wearable}. Workplace interventions such as Take-a-Stand and Move More @ Work focus on changing sitting patterns through organisational or timed prompts \cite{pronk2012take,hargreaves2023move}. In contrast, this thesis focuses on whether the relevant pelvic movement can be sensed and acted on locally by a body-worn research device. It therefore does not replace behavioural intervention studies; it supplies a technical component that could later be evaluated inside such a study.

The practical value of the prototype is its integration. Each individual part already exists in related work. IMUs can measure spine and lower-trunk movement in wearable assessment settings \cite{papi2017wearable}. Feature-based classifiers have been used for real-time activity recognition from inertial signals \cite{khan2013exploratory}. Haptic or vibrotactile feedback has also been used in posture-related wearable systems \cite{gopalai2011vibrotactile_posture_system}. The contribution here is that these parts are connected in one pelvis-focused pipeline and tested with actual field recordings. This is also why the work should be judged as a research prototype and feasibility study rather than as a polished engineering product or validated intervention.

The body-worn design also differs from camera and furniture-based approaches. Vision systems can classify sitting posture in prepared spaces \cite{estrada2023machine}, and smart chairs or cushions can measure pressure patterns at one seat \cite{tan2001sensing,ma2017activity}. The present system sacrifices full posture reconstruction and seat-level pressure detail in order to follow the user across seats. That trade-off fits the thesis aim, because the target is not clinical posture diagnosis but detection of whether pelvic movement has occurred during prolonged sitting.

The mechanical result also shows the current engineering boundary. The prototype was sufficient for supervised field data collection and short wearing sessions, but it is not yet a robust daily-wear device. The visible electronics and accessible wiring supported debugging and iteration, while the sterile-dressing attachment gave a simple short-session wearing method. These same choices limit long-term robustness. The intermittent battery-contact gaps observed during field collection show that the next hardware iteration should prioritise power-contact reliability, enclosure integration, and more repeatable body attachment.

\section{Objective 2: Detection Performance and Generalisation Boundary}

The classification results are strong for the labelled pilot task but narrow in behavioural meaning. The task separates static sitting from instructed anterior-posterior pelvic rotation. It is therefore simpler than broader sensor-based activity recognition, where systems often distinguish many daily activities \cite{trabelsi2022sensor}, and different from sitting-posture classification, where the target may be a full posture category rather than the presence of a movement event \cite{kappattanavar2021monitoring}. The high KNN F1-score shows that the two instructed classes are separable in the collected feature space, not that all spontaneous seated movements will be recognised.

The classifier comparison also explains why F1-score is more informative than accuracy here. The held-out set is highly imbalanced, so a majority-class predictor would achieve high accuracy while detecting no pelvic rotation. KNN gave the best balance between precision and recall. For the reminder logic, false positives are important because they would reset the inactivity timer when the user had not moved; false negatives are also relevant because they could fail to reset the timer after a real pelvic movement. The KNN result contains both kinds of errors, but the false-positive count is especially low. SVM produced higher recall but much lower precision, which would make the reminder less reliable, while random forest was more conservative and missed more pelvic-rotation windows. Logistic regression performed poorly because the separation is unlikely to be well represented by a single linear boundary in the 90-dimensional feature space.

The largest open question is naturalness. Participants performed an instructed pelvic tilt, whereas everyday seated movement is likely smaller, less periodic, and mixed with incidental shifts, reaching, phone use, or train motion. The classifier has therefore been tested as a controlled detector of an instructed movement pattern collected in train-sitting segments. A future dataset should include spontaneous seated behaviour in office, home, and transport contexts before the model can be described as robust for daily use.

\section{Objective 3: Sensor Validity, Placement, and Timing Limitations}

The Shimmer comparison supports trend-level use of the XIAO data for this detection task, but it should not be read as anatomical angle validation. The co-located recording showed that the custom node captured the timing of movement peaks when both devices were moved together, and the pilot field-worn subset showed similar long-term trends during the outdoor and train recordings. The NRMSE values indicate that the two devices did not produce identical amplitudes. This is consistent with the broader wearable-spine literature, where IMUs are useful for movement assessment but sensor position, calibration, and analysis goal strongly affect interpretation \cite{papi2017wearable}. For this thesis, amplitude disagreement is less critical than event timing because the classifier is trained and tested entirely on XIAO data.

The sacrum/lower-back placement follows the same logic as pelvis-focused IMU work by Liao et al., where sensor placement is treated as a central determinant of pelvic-rotation detection \cite{liao2026imu_pelvic_rotation}. The difference is scope. Liao et al. provide a benchmark-style dataset and placement comparison, whereas this thesis adds a low-cost prototype, mobile logging, haptic output, and embedded deployment. The cost of that integration is a smaller and less varied dataset.

The 1.2--1.3~s lag observed in the pilot field-worn recordings shows that real deployments are affected by the logging chain, not only by the sensor. Because the co-located lag was near zero, the field offset is most plausibly a clock and timestamping-chain issue, especially the independent Shimmer logger clock relative to the Android receive timestamps used for the XIAO log, rather than a delay of the inertial sensors themselves.

This distinction affects the strength of the validity claim. The Shimmer results make it reasonable to use the XIAO for a binary movement-detection task, but they do not establish that the wearable can estimate clinical pelvic kinematics. A stronger biomechanical validation would require controlled placement, a larger paired dataset, and comparison against a measurement setup designed for anatomical angles. For the present thesis, the relevant question is narrower: whether the low-cost device captures enough motion information to support the classifier.

\section{Objective 4: Embedded Reminder Feasibility and Haptic Boundary}

The embedded result connects the sensing task to TinyML constraints. Many sensor-based recognition studies use deep learning or relatively large offline models, but wearable microcontrollers impose limits on memory, computation time, and energy \cite{diab2022embedded,abadade2023tinyml}. The deployed model follows a conservative route: time-domain features, KNN, and KMeans prototype compression. KNN is normally a memory-heavy lazy learner because it stores the training set \cite{bonaccorso2017machine}. Prototype reduction and prototype generation are established ways to reduce the reference set used by nearest-neighbour classification \cite{hart1968condensed,triguero2012prototype_generation}. Replacing 9{,}391 windows with 128 prototypes changes the storage requirement and makes the method feasible on the XIAO nRF54L15 Sense.

The compact model should not be interpreted as strictly better than the full KNN baseline, even though its held-out F1-score is slightly higher, because both the feature set and the reference set changed. The more defensible interpretation is that prototype compression preserved the relevant decision behaviour while reducing the stored reference set by roughly 73 times. The held-out F1-score comes from the offline 2~s-window, 1~s-step evaluation, while the embedded timing measurement uses the deployed non-overlapping 2~s window cycle. The 2.73~ms per-window runtime is far below the 2~s window period, but power consumption was not measured. Timing feasibility therefore does not yet establish battery-life feasibility.

The DRV2605L gives the prototype a feedback channel that does not depend on a phone notification. That matters because the intended use case is a local cue during sitting, not post-hoc visualisation. However, the haptic part of the system has only been technically demonstrated. Sedentary-break studies evaluate whether interruptions to sitting change behavioural or physiological outcomes \cite{buffey2022_interrupting_sitting_meta}. Biofeedback studies examine whether users perceive, accept, and respond to feedback, including haptic or vibrotactile cues \cite{battis2024wearable}. This thesis does not yet answer those questions. The current 10~min inactivity threshold and 1~s vibration pulse are implementation choices, not validated behavioural parameters.

The complete system should therefore be described as a sensing, classification, and cue-delivery platform. It can detect the instructed movement in the pilot data and trigger a local reminder, but it cannot yet be claimed to reduce sedentary time, improve posture, or change health outcomes. Those claims would require a user study with measured reminder awareness, acceptance over repeated use, actual movement response, and longer-term adherence.

\section{Overall Limitations}

The limitations should be interpreted at four levels: dataset, analysis pipeline, hardware, and behavioural scope. At the dataset level, the classifier was trained on a small number of participants and tested on a single held-out participant. At the analysis level, clean seated segments and compact features were selected manually, although the final held-out participant was kept out of feature selection, scaling, and model-selection steps. At the hardware level, the prototype was sufficient for supervised field data collection but not yet robust enough for unsupervised long-term wear; the sterile-dressing attachment was stable for short pilot sessions but should not be treated as a final daily-wear solution. At the behavioural level, the system can deliver a cue, but the study did not test whether that cue changes sitting behaviour.

\begin{enumerate}
    \item \textbf{Small sample size.} Five participants, with one held out for testing. Population-level conclusions cannot be drawn from a single unseen participant.
    \item \textbf{No demographic stratification.} Age, sex, body composition, sitting habits, and low-back-pain history were not analysed, so inter-person differences cannot be explained clinically.
    \item \textbf{Instructed movement.} Pelvic rotation was an instructed anterior-posterior tilt, not spontaneous seated behaviour.
    \item \textbf{Single movement direction.} Lateral pelvic rotation and other seated mobility patterns were not labelled.
    \item \textbf{Manual segment and feature selection.} Clean segments and compact features were selected manually. The final participant was held out from selection, but the workflow should still be treated as preliminary rather than fully automated.
    \item \textbf{Limited Shimmer coverage.} Three of five paired Shimmer recordings from the pilot field study failed; field-worn reliability conclusions are based on two participants.
    \item \textbf{Context specificity.} All classifier data are from sitting during train journeys; office work and home use have not been tested.
    \item \textbf{Hardware reliability.} Intermittent battery-contact gaps could cause false inactivity detections in deployment.
    \item \textbf{No battery-life measurement.} On-device timing was measured, but power consumption and runtime under continuous sensing and feedback were not measured.
    \item \textbf{No user study.} The vibration reminder was actuated on-device but not evaluated for awareness, acceptability, or behavioural response.
\end{enumerate}

\section{Future Work}

A larger and more diverse cohort would support cross-validation with multiple held-out subjects and allow factors such as age, body mass index, and history of lower-back pain to be examined. The labelling protocol should be extended beyond instructed pelvic rotation to include spontaneous seated movement, lateral rotation, forward lean, and weight shifts. The hardware needs a more reliable battery connection and body attachment, and a comparison of sacrum, lower back, and waist positions would clarify the optimal placement. On the algorithm side, personalised calibration with a short user-specific recording could adapt the prototype set per user, and power profiling on the deployed firmware should clarify battery life under continuous IMU sampling and classification. A dedicated user study is needed before the system can be described as a behavioural intervention rather than a sensing and classification platform.

The next study should also separate technical validation from behavioural evaluation. A technical study should first test sensor reliability, battery life, placement stability, and classifier generalisation in office, home, and transport contexts. A later user study can then examine reminder timing, vibration intensity, annoyance, compliance, and actual movement response. Keeping those phases separate would make the evidence clearer: first show that the system detects the target behaviour reliably, then test whether the cue changes behaviour in a useful way.

A final practical step is to improve reproducibility of the embedded pipeline. The training scripts, selected features, prototype values, firmware header generation, and timing test should be kept as one documented workflow. That would make it easier to update the model when more participant data become available and reduce the risk that the offline and embedded versions diverge. In addition, future reports should state the exact firmware version, phone model, Android version, and BLE logging settings, because the field lag observed in this thesis shows that the data path can affect timing results.
